% ECONOMICS_PROBLEM: {"id": "D31-330", "title": "Productivity versus rents in top incomes", "jel_code": "D31", "source_id": "OP-0608"}
% !TeX program = xelatex
\documentclass[11pt,a4paper]{article}
\usepackage[margin=23mm]{geometry}
\usepackage{fontspec}
\setmainfont{Latin Modern Roman}
\usepackage{amsmath,amssymb,mathtools}
\usepackage{xcolor,enumitem,booktabs,longtable,array,xurl,hyperref}
\definecolor{muted}{HTML}{53636C}
\hypersetup{colorlinks=true,urlcolor=blue,linkcolor=blue}
\setlength{\parindent}{0pt}
\setlength{\parskip}{5pt}
\newcommand{\R}{\mathbb R}
\newcommand{\E}{\mathbb E}
\newcommand{\N}{\mathbb N}
\newcommand{\ind}{\mathbf 1}
\newcommand{\Var}{\operatorname{Var}}
\newcommand{\Cov}{\operatorname{Cov}}
\newcommand{\supp}{\operatorname{supp}}
\DeclareMathOperator*{\argmin}{arg\,min}
\DeclareMathOperator*{\argmax}{arg\,max}
\newcommand{\entrymeta}[3]{{\sffamily\small\color{muted}\textbf{\texttt{#1}}\quad|\quad #2\par #3\par}}
\newcommand{\Problemref}[1]{\texttt{#1}}
\newcommand{\JELref}[2]{\texttt{#2} (JEL \texttt{#1})}
\begin{document}
\section*{Productivity versus rents in top incomes}
\textbf{Problem ID:} \texttt{D31-330}\quad \textbf{JEL:} \texttt{D31}\par
% BEGIN ECONOMICS STATEMENT
\label{problem:OP-0608}
% ECONOMICS_PROBLEM: {"local_id": "ECON-029", "title": "Productivity versus rents in top incomes", "kind": "R", "audit_date": "2026-10-08", "source_identity_checked": true, "agenda_or_question_support_checked": true, "source_role": "see per-entry audit and locators", "current_open_certified": false}


\label{audited:ECON-029}

{\sffamily\small\color{muted}\textbf{ECON-029} | R | Source-backed agenda; exact formalization is editorial\par}

\subsection*{Definitions and formal target}

In a specified production and compensation model write \(y_i=\operatorname{MP}_i+\rho_i\), where marginal social output \(\operatorname{MP}_i\) uses the same currency and horizon as pay. For \(q=F_y^{-1}(1-p)\), let \(\phi_p(y)\) be one above \(q\), zero below and a fraction at \(q\) chosen so \(\mathbb E[\phi_p(y_i)]=p\). For finite moments and positive denominator, characterize the identified set for
\[
 \theta_\rho=\frac{\mathbb E[\rho_i\phi_p(y_i)]}
                   {\mathbb E[y_i\phi_p(y_i)]}.
\]
Data may include pay, firm outcomes, worker moves and exogenous governance changes.

\subsection*{Clarifications and required assumptions}

Private marginal revenue product differs from social marginal product when there are rents or externalities. Specify risk compensation, human-capital investment, deferred pay and outside options before interpreting the residual as rent. It can be negative; nonnegative rents need an extra restriction. Define output, population and compensation components. High pay and firm performance alone do not identify marginal contribution. Fractional ranking makes the top-\(p\) group exact with atoms.

\subsection*{Variants and changes of scope}

\begin{itemize}
\item Executives, entrepreneurs and employees require different compensation models.
\item A residual differs from the welfare effect of eliminating rent-generating institutions, since behavior can respond.
\item Fixed baseline ranks differ from counterfactual re-ranking.
\end{itemize}

\subsection*{Audit conclusion and source relationship}

[S1], PDF p. 7, supports productivity versus rents in top incomes. This ratio is editorial and does not classify every difference from a productivity proxy as rent.

\subsection*{References and exact source roles}

{\footnotesize\raggedright
\par\noindent\textbf{[S1]} Benjamin Moll (n.d.). \emph{Lecture 11: Good Research Topics and Open Questions}. Distributional Macroeconomics lecture slides; publicly archived at a 2019 URL.\par \textit{Verified locator / source role:} PDF p. 7: top incomes, productivity, and rent seeking.\par \url{https://benjaminmoll.com/wp-content/uploads/2019/07/Lecture11_2149.pdf}\par\smallskip
}

\subsection*{Common conventions: Additional-source status and mathematical conventions}

\label{audited:conventions}

Of the 100 supplied ECON entries, 65 distinct problems appear here; 32 close
matches refine earlier entries, and three duplicate questions map to existing
entries. Appendix~\ref{app:audited-crosswalk} lists every destination.
The attachment's P/R/S/U distinctions and source audits are retained. Its
precise mathematical formalizations are editorial unless the individual audit
says otherwise. Candidate subsidy targets ECON-004--006 retain their U flags;
the resolved approval-core question ECON-007 updates OP-0202 above.
The following supplied conventions govern both this part and the attributed
ECON refinements elsewhere in the catalogue, subject to local restrictions.

\textbf{Parameterized questions.} A problem family lists inputs that must be fixed before an instance is evaluated: population, horizon, policies, information, data law, model class and welfare weights. These are required choices, not quantities the citations are assumed to specify. Undefined applications should not be treated as identified estimands.

\textbf{Indices and integrability.} Sets of agents, firms and regions are finite unless a population measure is explicitly used. \(i,j\) index units, \(r\) regions, \(t\) dates and \(h\) horizons. \(\mathbb E\) and \(\Pr\) refer to the declared population and law; \(\mathbf1\{A\}\) is an indicator. Every displayed expectation and discounted sum needs existence in the stated domain. Infinite horizons use a discount in \((0,1)\) and a convergence condition; finite horizons may use discount one.

\textbf{Optimization and existence.} A displayed \(\max\), \(\min\), or \(\arg\max\) is conditional on attainment. For finite feasible menus this is immediate; general spaces need continuity and compactness/coercivity or another existence argument. Without attainment, the target is the supremum/infimum and arbitrarily near-optimal feasible choices. \(\inf\varnothing=+\infty\). Empty feasibility and an unidentified target are different issues.

\textbf{Potential outcomes.} \(Y_i(z)\) is the outcome under a specified intervention, not the observed conditional mean among people choosing \(z\). In binary comparisons \(\Delta Y_i=Y_i(1)-Y_i(0)\). Specify assignment, treatment versions, compliance, information and observation horizon. Interference is allowed under a full policy vector; compressed exposure notation needs a justified mapping. No interference, consistency and overlap are substantive assumptions, not automatic consequences of notation.

\textbf{Populations and observation.} Fix baseline eligibility and population weights. Conditioning on post-policy employment, survival, relocation or program uptake can change the population being compared. Death is not ordinary loss to follow-up; outcomes truncated by death need a composite convention, joint survival target, or explicitly defined survivor stratum. Fertility-changing interventions need a population functional rather than an unexplained fixed descendant cohort. Measurement and missingness models must be supplied.

\textbf{Identification.} A structural model \(m\in\mathcal M\) implies observable law \(P_m\) and target \(\theta(m)\). Identification means
\[
 P_{m_1}=P_{m_2}\ \Longrightarrow\ \theta(m_1)=\theta(m_2)
 \quad\text{for all }m_1,m_2\in\mathcal M .
\]
Without identification, the target is
\[
 \Theta(P;\mathcal M)=\{\theta(m):m\in\mathcal M,\ P_m=P\}.
\]
Writing \(\theta(P)\) before establishing identification can hide incompatible latent models. Confidence sets additionally need a sampling law, confidence level, asymptotic sequence and uniformity class. Point identification, coverage and informative precision are separate properties.

\textbf{Mediation.} Total effects of randomized assignment are distinct from cross-world natural indirect effects. Belief/effort-channel effects require a meaningful mediator intervention and additional measurement, exclusion or mediation assumptions. Randomization of the initial policy alone does not supply those assumptions. Report bounds or interventional alternatives when the stronger target is unsupported.

\textbf{Equilibrium.} A counterfactual \(W(z)\), \(p(z)\), or \(\mu_t^z\) requires individual optimization, feasibility, government/resource budgets, market clearing and state-transition laws. State equilibrium existence and selection, or report ranges over compatible equilibria. Initial-state integration includes subsequent endogenous transitions; current-state integration must use the policy-dependent distribution. Reduced-form invariance after policy is not assumed.

\textbf{Welfare accounts.} Scores, utility, health, dollars and ecological quantities require explicit conversion before addition. Fees, taxes, subsidies, wages and extortion payments are transfers between parties; distributional weights can make them welfare relevant, but their face value is not automatically a resource loss. State ownership, geographic boundaries, foreign-income weights and fiscal finance. Do not add profits to household welfare already including dividends, or count land capitalization and the underlying benefit twice. A benefit-cost expression must distinguish gross benefits from net welfare and subtract every real cost exactly once.

\textbf{Derivatives and risk.} Log derivatives require positive arguments; local responses require admissible perturbations and specified equilibrium branches. Differentiation under expectations needs regularity/domination, and boundaries may allow only one-sided derivatives. Risk uses a specified law; ambiguity uses a declared nonempty law set on a common history space. Adaptive policies must be nonanticipative. A robust criterion is an editorial choice unless attributed explicitly.

\textbf{Scope overlaps.} Allocation, collective choice, contracts, household bargaining and coordinated policy overlap economics with optimization and game theory. They are retained because they appeared in the original catalogue; no claim of a strictly game-theory-free selection is made.
% END ECONOMICS STATEMENT
\end{document}
